\paraid{p-e72a35de4049}
Let
$\GHSpace$
denote the space of isometry classes of nonempty compact metric spaces
with the Gromov--Hausdorff distance
$\GHDistance$.
In this paper,
we construct a measure for each compact metric space.
The assignment is continuous in the Gromov--Hausdorff topology.
Specifically,
if the spaces are isometrically embedded in a common compact metric space
and converge in Hausdorff distance,
then the pushed-forward measures converge weakly.

\paraid{p-96038747b89a}
Larrieu
\cite[arXiv v4, Theorem~3]{Larrieu2014InvariantProbabilities}
describes a construction of Borel probability measures on nonempty compact
metric spaces that are invariant under isometries between open subsets.
The construction uses a fixed generalized limit.
The cited result does not assert full support or continuity as the underlying
compact metric space varies.
Leinster and Roff
\cite[Theorem~1.1 and Definition~9.2]{LeinsterRoff2021Entropy}
study probability measures that maximize their family of metric entropies.
They define a uniform measure by taking a weak limit of maximizing measures
as the metric is multiplied by a factor tending to infinity,
under the assumptions that the maximizing measure is unique for all sufficiently
large factors and that the limit exists.
The  problem here concerns the ordinary Gromov--Hausdorff topology
on the underlying compact spaces and weak convergence of measures in
common ambient spaces.

\paraid{p-6e73278858c1}
In our main result,
\Cref{thm:invariant-measures},
we assign to every nonempty compact metric space
$\MetricSpace{\BaseCarrier}{\MetricSymbol}$
a Borel probability measure
$\SelectedMeasure{\BaseCarrier}{\MetricSymbol}$
with full support,
so that
\[
 \supp(\SelectedMeasure{\BaseCarrier}{\MetricSymbol})=\BaseCarrier.
\]
For every isometry
$\IdentifyingIsometry\colon\MetricSpace{\BaseCarrier}{\MetricSymbol}
\to\MetricSpace{\ComparisonCarrier}{\ComparisonMetric}$,
we have
\[
 \IdentifyingIsometry\Pushforward\SelectedMeasure{\BaseCarrier}{\MetricSymbol}
 =\SelectedMeasure{\ComparisonCarrier}{\ComparisonMetric}.
\]
This means well-definedness.
Let compact subspaces
$\BaseCarrier_\SequenceIndex$
and
$\BaseCarrier$
of a compact metric space
$\MetricSpace{\AmbientCarrier}{\AmbientMetric}$
carry the restricted metrics
$\MetricSymbol_\SequenceIndex$
and
$\MetricSymbol$.
Then
\[
 \HausdorffDistance{\AmbientMetric}(\BaseCarrier_\SequenceIndex,\BaseCarrier)\to0
 \quad\Longrightarrow\quad
 \SelectedMeasure{\BaseCarrier_\SequenceIndex}{\MetricSymbol_\SequenceIndex}
 \to\SelectedMeasure{\BaseCarrier}{\MetricSymbol}
 \quad\text{weakly in }\ProbabilityMeasures(\AmbientCarrier).
\]
We regard the selected measures as probabilities on the common ambient space.
Since
$\supp(\SelectedMeasure{\BaseCarrier}{\MetricSymbol})=\BaseCarrier$,
for every
$1\leq\IntegrabilityExponent<\infty$
and every
$\HilbertFunction\in\ContinuousFunctions(\BaseCarrier)\setminus\{0\}$,
we have
\[
 \norm{\HilbertFunction}_{\LebesgueSpace^\IntegrabilityExponent
   (\BaseCarrier,\SelectedMeasure{\BaseCarrier}{\MetricSymbol})}^{\IntegrabilityExponent}
 =\int_\BaseCarrier|\HilbertFunction(\BasePoint)|^{\IntegrabilityExponent}
   \,\IntegrationDifferential\SelectedMeasure{\BaseCarrier}{\MetricSymbol}(\BasePoint)
 >0.
\]
Isometry compatibility and weak continuity permit the corresponding
integrals to be compared as the compact spaces vary.

\paraid{p-a9e34dfe60e6}
To obtain an invariant full-support probability,
for a fixed compact metric space,
we average a full-support probability over its compact isometry group.
Rouyer's results
\cite[arXiv v1, Theorems~2 and~4]{Rouyer2011}
imply that compact metric spaces with trivial isometry group are dense in
$\GHSpace$.
Thus the map assigning to each compact metric space its isometry group,
equipped with the uniform metric, is discontinuous with respect to the
Gromov--Hausdorff distance.
We therefore construct the invariant probabilities simultaneously and continuously.
We work in the Gromov--Hausdorff--Prokhorov space of compact spaces
with probability measures.
In this space,
we retain the whole underlying compact space even when
the measure has smaller support.
Let
$\InvariantMeasuredSpaces$
denote the subspace of invariant full-support measured spaces.
We prove that
$\InvariantMeasuredSpaces$
is Polish and that the forgetful map,
\[
 \MeasureProjection\colon\InvariantMeasuredSpaces\to\GHSpace,
 \qquad
 \MeasureProjection\MeasuredSpace{\BaseCarrier}{\MetricSymbol}{\FirstMeasure}
 =\MetricSpace{\BaseCarrier}{\MetricSymbol},
\]
is a continuous open surjection.
By the theorems of Banakh and Radul
\cite[Theorem~1.2, p.~20]{BanakhRadul1999}
and Valov (\Cref{thm:valov}),
there is a continuous map
\[
 \LawSection\colon\GHSpace\to\ProbabilityMeasures(\InvariantMeasuredSpaces),
\]
satisfying
\[
 \MeasureProjection\Pushforward\LawSection\MetricSpace{\BaseCarrier}{\MetricSymbol}
 =\DiracProbability_{\MetricSpace{\BaseCarrier}{\MetricSymbol}}.
\]
Fix a compact metric space
$\MetricSpace{\BaseCarrier}{\MetricSymbol}$.
Each point of the fiber
$\MeasureProjection^{-1}(\MetricSpace{\BaseCarrier}{\MetricSymbol})$
can be represented by a compact metric space isometric to
$\MetricSpace{\BaseCarrier}{\MetricSymbol}$,
together with an invariant full-support probability measure.
We choose an isometry from that space to
$\MetricSpace{\BaseCarrier}{\MetricSymbol}$
and push the measure forward to
$\BaseCarrier$.
The resulting probability does not depend on the chosen isometry,
since the measure is invariant under isometries.
It also does not depend on the representative,
since two representatives of the same class are related by a measure-preserving isometry.
We average these probabilities with respect to
$\LawSection\MetricSpace{\BaseCarrier}{\MetricSymbol}$,
using the Riesz--Markov--Kakutani theorem
(\Cref{thm:riesz-markov-kakutani}).
We prove continuity of the selected measures by comparing integrals in a
common compact ambient space.

\paraid{p-b5920f7184fb}
The law and its average have different domains.
For the fixed representative
$\MetricSpace{\BaseCarrier}{\MetricSymbol}$,
put
$\MeasureFiber{\BaseCarrier}
 =\MeasureProjection^{-1}(\MetricSpace{\BaseCarrier}{\MetricSymbol})$.
For every
$\FiberPoint\in\MeasureFiber{\BaseCarrier}$,
let
$\FiberMeasure{\BaseCarrier}{\FiberPoint}$
be the transported probability on
$\BaseCarrier$
described above.
For every continuous real function
$\TestFunction$
on
$\BaseCarrier$,
the selected measure satisfies
\begin{equation}\label{eq:part-i-average-interface}
 \int_\BaseCarrier\TestFunction(\BasePoint)
 \,\IntegrationDifferential\SelectedMeasure{\BaseCarrier}{\MetricSymbol}(\BasePoint)
 =\int_{\MeasureFiber{\BaseCarrier}}
 \left(\int_\BaseCarrier\TestFunction(\BasePoint)
 \,\IntegrationDifferential\FiberMeasure{\BaseCarrier}{\FiberPoint}(\BasePoint)\right)
 \,\IntegrationDifferential\LawSection\MetricSpace{\BaseCarrier}{\MetricSymbol}(\FiberPoint).
\end{equation}
The inner integral uses an invariant full-support probability on the
fixed carrier,
and the outer integral uses a law on measured isomorphism classes.
\Cref{lem:continuous-barycenter} shows that this averaging preserves
invariance and full support.

\paraid{p-3423d1d0cc93}
As an application,
we construct a topological embedding
$s$
of
$\GHSpace$
into the Gromov--Hausdorff--Prokhorov space by attaching the selected measure.
The forgetful map is a continuous left inverse of
$s$.
We prove that the image
$s(\GHSpace)$
is a retract of the Gromov--Hausdorff--Prokhorov space and is contained in
the subspace of invariant full-support probabilities
(Corollary~\ref{cor:ghp-retract}).

\paraid{p-7734020dafc8}
For metric measure spaces,
Kazukawa,
Nakajima,
and Shioya
\cite[arXiv v1, Theorems~1.4--1.6]{KazukawaNakajimaShioya2024}
prove contractibility and local path connectedness for both the box and
concentration topologies.
The section and retraction constructed here relate the ordinary
Gromov--Hausdorff space to the Gromov--Hausdorff--Prokhorov space.

\medskip
\noindent\textbf{Dependence on the other parts.}
\paraid{p-f205aa0a2b93}
We prove the selection theorem and the retraction using the common
preliminaries and the results recalled in Section~\ref{sec:prelim-measures}.
In Part~\ref{part:spectral},
we use \Cref{thm:invariant-measures}
to construct local models.
The later absolute retract argument combines these local models.
The selected measures enter that argument in the construction of the local models.

\medskip
\noindent\textbf{Organization.}
\paraid{p-0dd0534490f6}
Section~\ref{sec:prelim-measures} contains the preliminaries for this part.
In Subsection~\ref{subsec:measures-ghp},
we use the probability background from Subsection~\ref{subsec:preliminaries-probability}
to introduce the Gromov--Hausdorff--Prokhorov space and prove common ambient
realization results.
Subsection~\ref{subsec:measures-preliminaries-laws} recalls Haar probability
and Valov's theorem.
In Subsection~\ref{subsec:measures-open-projection},
we prove that the space of metric measured spaces with invariant full-support
measures is Polish
and that its projection onto the unmeasured Gromov--Hausdorff space is open.
In Subsection~\ref{subsec:measures-laws-and-averages},
we construct continuous probability laws on the fibers
and average them to prove \Cref{thm:invariant-measures}.
In Subsection~\ref{subsec:measures-retract},
we identify the selected measured spaces as a retract
of the Gromov--Hausdorff--Prokhorov space.

\medskip
\noindent\textbf{Conventions and notation.}
\paraid{p-33d1f4dd98d2}
All compact metric spaces are nonempty.
When a metric space represents a point of
$\GHSpace$,
we identify it with its isometry class.
An isometry is a surjective isometric embedding.
We denote by
$\ProbabilityMeasures(\DomainSpace)$
the set of Radon probability measures on a metrizable space
$\DomainSpace$,
equipped with the topology of weak convergence against bounded continuous real functions.
Subscripts in
$\SelectedMeasure{\BaseCarrier}{\MetricSymbol}$
refer to the specified metric space
$\MetricSpace{\BaseCarrier}{\MetricSymbol}$.
For isometrically embedded subsets of a common ambient space,
we use the same symbols for the subsets and their restricted metrics.
Sequences are indexed by
$\NonnegativeIntegers=\{0,1,2,\ldots\}$.
